% Test Scenarios for Chapter 5 - Performance Tests
% Copy this content into latex/tex/5-testy-wydajnosci.tex

\clearpage
\section{Testy wydajności}
\label{sec:testy-wydajnosci}

\subsection{Metodyka przeprowadzania testów}

\subsubsection{Przygotowanie środowiska testowego}

Wszystkie testy wydajnościowe przeprowadzono w~kontrolowanych warunkach, zapewniających powtarzalność i~porównywalność wyników. Przed każdym cyklem pomiarowym wykonano następujące kroki przygotowawcze:

\begin{enumerate}
    \item Zamknięto wszystkie aplikacje działające w~tle
    \item Wyłączono automatyczne aktualizacje systemu operacyjnego
    \item Ustawiono tryb zasilania na ,,Wysoka wydajność''
    \item Wyłączono oszczędzanie energii procesora i~karty graficznej
    \item Odczekano 5~minut na stabilizację temperatury komponentów
    \item Wykonano trzy pomiary próbne (,,rozgrzewające'') przed rejestracją właściwych wyników
\end{enumerate}

\subsubsection{Standaryzacja warunków testowych}

Aby zapewnić porównywalność wyników między Unity a~Unreal Engine, opracowano trzy scenariusze testowe o~zróżnicowanym poziomie obciążenia. Każdy scenariusz został zaimplementowany w~identyczny sposób w~obu silnikach, z~następującymi parametrami:

\begin{table}[h!]
    \centering
    \caption{Parametry scenariuszy testowych}
    \label{tab:test-scenarios}
    \begin{tabular}{|l|c|c|c|}
        \hline
        \textbf{Parametr} & \textbf{Niski} & \textbf{Średni} & \textbf{Wysoki} \\
        \hline\hline
        Liczba pocisków & 50--100 & 200--300 & 500+ \\
        \hline
        Liczba przeciwników & 2--3 & 5--7 & 10+ \\
        \hline
        Czas trwania testu & 30 s & 30 s & 30 s \\
        \hline
        Liczba pomiarów & 5 & 5 & 5 \\
        \hline
    \end{tabular}
\end{table}

\subsection{Scenariusze testowe gry bullet-hell}

\subsubsection{Scenariusz 1: Niski poziom trudności (baseline)}

\paragraph{Cel testu}
Ustalenie wydajności bazowej przy minimalnym obciążeniu systemu renderowania i~fizyki.

\paragraph{Parametry}
\begin{itemize}
    \item \textbf{Liczba pocisków na ekranie}: 50--100 jednocześnie
    \item \textbf{Aktywni przeciwnicy}: 2--3 jednostki
    \item \textbf{Czas trwania pomiaru}: 30 sekund
    \item \textbf{Liczba powtórzeń}: 5 przechwytów klatek w~odstępach 5-sekundowych
\end{itemize}

\paragraph{Oczekiwane rezultaty}
W~scenariuszu bazowym oczekiwano stabilnej częstotliwości odświeżania na poziomie 60~FPS, niskiego wykorzystania GPU~(<50\%) oraz minimalnego zużycia pamięci.

\subsubsection{Scenariusz 2: Średni poziom trudności}

\paragraph{Cel testu}
Ocena wydajności przy umiarkowanym obciążeniu systemu, symulująca typową rozgrywkę.

\paragraph{Parametry}
\begin{itemize}
    \item \textbf{Liczba pocisków na ekranie}: 200--300 jednocześnie
    \item \textbf{Aktywni przeciwnicy}: 5--7 jednostek
    \item \textbf{Czas trwania pomiaru}: 30 sekund
    \item \textbf{Liczba powtórzeń}: 5 przechwytów klatek w~odstępach 5-sekundowych
\end{itemize}

\paragraph{Oczekiwane rezultaty}
Przewidywano umiarkowane wykorzystanie GPU~(50--70\%), możliwe niewielkie spadki częstotliwości klatek oraz wzrost zużycia pamięci.

\subsubsection{Scenariusz 3: Wysoki poziom trudności (test obciążeniowy)}

\paragraph{Cel testu}
Weryfikacja wydajności w~ekstremalnych warunkach przy maksymalnym obciążeniu systemu.

\paragraph{Parametry}
\begin{itemize}
    \item \textbf{Liczba pocisków na ekranie}: 500+ jednocześnie
    \item \textbf{Aktywni przeciwnicy}: 10+ jednostek
    \item \textbf{Czas trwania pomiaru}: 30 sekund
    \item \textbf{Liczba powtórzeń}: 5 przechwytów klatek w~odstępach 5-sekundowych
\end{itemize}

\paragraph{Oczekiwane rezultaty}
W~scenariuszu obciążeniowym spodziewano się wysokiego wykorzystania GPU~(>70\%), potencjalnych spadków wydajności oraz maksymalnego zaobserwowanego zużycia pamięci.

\subsection{Metryki wydajności}

\subsubsection{Zbierane dane}

Dla każdego scenariusza i~silnika rejestrowano następujące metryki przy użyciu NVIDIA Nsight Graphics:

\begin{itemize}
    \item \textbf{Czas klatki} (frame time) -- czas renderowania pojedynczej klatki w~milisekundach
    \item \textbf{FPS} (frames per second) -- liczba klatek na sekundę, wyliczana jako $1000 / \text{frame time}$
    \item \textbf{Wykorzystanie GPU} -- procent wykorzystania mocy obliczeniowej karty graficznej
    \item \textbf{Zużycie pamięci VRAM} -- ilość zajętej pamięci karty graficznej w~megabajtach
    \item \textbf{Liczba wywołań rysowania} (draw calls) -- liczba instrukcji renderowania na klatkę
    \item \textbf{Liczba wierzchołków} -- całkowita liczba przetworzonych wierzchołków na klatkę
\end{itemize}

\subsection{Wyniki testów}

\subsubsection{Tabele z wynikami pomiarów}

% Placeholder table - fill with actual data
\begin{table}[h!]
    \centering
    \caption{Wyniki testów wydajności -- Scenariusz niski}
    \label{tab:results-low}
    \begin{tabular}{|l|c|c|c|c|}
        \hline
        \textbf{Metryka} & \textbf{Unity śr.} & \textbf{Unity odch.} & \textbf{Unreal śr.} & \textbf{Unreal odch.} \\
        \hline\hline
        Czas klatki [ms] & [DATA] & [DATA] & [DATA] & [DATA] \\
        \hline
        FPS & [DATA] & [DATA] & [DATA] & [DATA] \\
        \hline
        GPU [\%] & [DATA] & [DATA] & [DATA] & [DATA] \\
        \hline
        VRAM [MB] & [DATA] & [DATA] & [DATA] & [DATA] \\
        \hline
        Draw calls & [DATA] & [DATA] & [DATA] & [DATA] \\
        \hline
        Wierzchołki [tys.] & [DATA] & [DATA] & [DATA] & [DATA] \\
        \hline
    \end{tabular}
\end{table}

\begin{table}[h!]
    \centering
    \caption{Wyniki testów wydajności -- Scenariusz średni}
    \label{tab:results-medium}
    \begin{tabular}{|l|c|c|c|c|}
        \hline
        \textbf{Metryka} & \textbf{Unity śr.} & \textbf{Unity odch.} & \textbf{Unreal śr.} & \textbf{Unreal odch.} \\
        \hline\hline
        Czas klatki [ms] & [DATA] & [DATA] & [DATA] & [DATA] \\
        \hline
        FPS & [DATA] & [DATA] & [DATA] & [DATA] \\
        \hline
        GPU [\%] & [DATA] & [DATA] & [DATA] & [DATA] \\
        \hline
        VRAM [MB] & [DATA] & [DATA] & [DATA] & [DATA] \\
        \hline
        Draw calls & [DATA] & [DATA] & [DATA] & [DATA] \\
        \hline
        Wierzchołki [tys.] & [DATA] & [DATA] & [DATA] & [DATA] \\
        \hline
    \end{tabular}
\end{table}

\begin{table}[h!]
    \centering
    \caption{Wyniki testów wydajności -- Scenariusz wysoki}
    \label{tab:results-high}
    \begin{tabular}{|l|c|c|c|c|}
        \hline
        \textbf{Metryka} & \textbf{Unity śr.} & \textbf{Unity odch.} & \textbf{Unreal śr.} & \textbf{Unreal odch.} \\
        \hline\hline
        Czas klatki [ms] & [DATA] & [DATA] & [DATA] & [DATA] \\
        \hline
        FPS & [DATA] & [DATA] & [DATA] & [DATA] \\
        \hline
        GPU [\%] & [DATA] & [DATA] & [DATA] & [DATA] \\
        \hline
        VRAM [MB] & [DATA] & [DATA] & [DATA] & [DATA] \\
        \hline
        Draw calls & [DATA] & [DATA] & [DATA] & [DATA] \\
        \hline
        Wierzchołki [tys.] & [DATA] & [DATA] & [DATA] & [DATA] \\
        \hline
    \end{tabular}
\end{table}

\subsubsection{Placeholder dla wykresów}

% Placeholder - replace with actual figures generated from data
\begin{figure}[h!]
    \centering
    % \includegraphics[width=0.8\linewidth]{tex/img/frame-time-comparison.pdf}
    \caption{Porównanie czasu klatki między Unity a~Unreal Engine w~trzech scenariuszach testowych. [PLACEHOLDER -- wygeneruj wykres za pomocą skryptu scripts/generate\_plots.py]}
    \label{fig:frame-time-comparison}
\end{figure}

\begin{figure}[h!]
    \centering
    % \includegraphics[width=0.8\linewidth]{tex/img/gpu-utilization.pdf}
    \caption{Wykorzystanie GPU w~funkcji liczby obiektów na ekranie. [PLACEHOLDER -- wygeneruj wykres za pomocą skryptu scripts/generate\_plots.py]}
    \label{fig:gpu-utilization}
\end{figure}

\begin{figure}[h!]
    \centering
    % \includegraphics[width=0.8\linewidth]{tex/img/memory-usage.pdf}
    \caption{Zużycie pamięci VRAM w~trzech scenariuszach testowych. [PLACEHOLDER -- wygeneruj wykres za pomocą skryptu scripts/generate\_plots.py]}
    \label{fig:memory-usage}
\end{figure}

\subsection{Analiza wyników}

\subsubsection{Wydajność w scenariuszu bazowym}

[ANALIZA -- wypełnij po zebraniu danych]

Jak przedstawiono w~Tabeli~\ref{tab:results-low}, w~scenariuszu o~niskim obciążeniu oba silniki...

\subsubsection{Wydajność w scenariuszu średnim}

[ANALIZA -- wypełnij po zebraniu danych]

Wyniki przedstawione w~Tabeli~\ref{tab:results-medium} wskazują, że...

\subsubsection{Wydajność w scenariuszu obciążeniowym}

[ANALIZA -- wypełnij po zebraniu danych]

Najbardziej wymagający scenariusz (Tabela~\ref{tab:results-high}) ujawnił...

\subsection{Podsumowanie wyników testów wydajności}

\begin{table}[h!]
    \centering
    \caption{Zestawienie zbiorcze wyników testów}
    \label{tab:summary}
    \begin{tabular}{|l|c|c|c|}
        \hline
        \textbf{Kryterium} & \textbf{Unity} & \textbf{Unreal} & \textbf{Zwycięzca} \\
        \hline\hline
        Średni FPS (niski) & [DATA] & [DATA] & [TBD] \\
        \hline
        Średni FPS (średni) & [DATA] & [DATA] & [TBD] \\
        \hline
        Średni FPS (wysoki) & [DATA] & [DATA] & [TBD] \\
        \hline
        Zużycie VRAM & [DATA] & [DATA] & [TBD] \\
        \hline
        Efektywność draw calls & [DATA] & [DATA] & [TBD] \\
        \hline
    \end{tabular}
\end{table}
